All clauses are exact rational-function identities, verified at two masses, the angular clauses at the equator; general position is literally the same algebra, since $\sin\theta$ enters $\mu_2$ and $e_{\hat\phi}$ reciprocally and cancels identically in every pairing and Gram entry used (Section \ref{sec:verification}). (1) is the pairing table: $\dtPG(V) = 1$, $(k/\sqrt r)(e_{\hat r}) = 1$, $\mu_i(e_{\hat\imath}) = 1$, all others zero. (2): $G(\dtPG,\dtPG) = 1$, $G(\dtPG, k/\sqrt r) = 0$, $G(k/\sqrt r, k/\sqrt r) = g_k/r = -1$ exactly, by Theorem \claimref{sheaf.schwarzschild-degeneration-boundary-tame}. (3): membership in $\shW$ gives bounded components by (1) applied to the frame. Conversely, write $\alpha = A\,du + B\,dt + C'\,d\theta + D'\sin\theta\,d\phi$ with $u$-smooth coefficients: boundedness of $\alpha(e_{\hat r}) = A/(2u)$ forces $A = O(u)$, then boundedness of $\alpha(V) = B - mA/(2u^2)$ forces $A = O(u^2)$, and boundedness of the angular components $C'/u^2$, $D'/u^2$ forces $C', D' = O(u^2)$; a $u$-smooth function that is $O(u^2)$ is $u^2$ times a $u$-smooth function (Taylor division), so $\alpha$ is a $C^\infty(\Mt_u)$-combination of $\rho$, $dt$, $\mu_1$, $\mu_2$: $\alpha \in \shW$. The norm identity itself is the statement that the frame of (1) is $G$-orthonormal, i.e.\ clause (2). (4): the displayed transition and its determinant $\det\begin{psmallmatrix} -m & \sqrt r\\ 1 & 0\end{psmallmatrix} = -\sqrt r$; the frame comparison with Proposition \claimref{sheaf.exceptional-algebra-black-hole} replaces the timelike leg $\rho$ by $\dtPG$ through the same transition.
